%=====================================================================
\chapter{Business Models and Revenue Models}\label{ch:models}
%=====================================================================
\locator{1}{Part I --- E-Commerce and Its Models}

\begin{outcomes}
By the end of this chapter you will be able to:
\begin{tight}
  \item \textbf{Distinguish} a business model from a revenue model, and
        \textbf{explain} why a venture can have a sound one and a hopeless
        other.
  \item \textbf{Classify} an online business by participant --- B2C, B2B, C2C,
        C2B --- and say what each classification predicts about the software.
  \item \textbf{Name} the principal revenue models and \textbf{identify} which
        one a given shop is using.
  \item \textbf{Explain} the unit economics of a single order, and
        \textbf{compute} the contribution margin from a set of figures.
  \item \textbf{Argue} which revenue model a given venture should adopt, and
        \textbf{state} the software consequences of that choice.
\end{tight}
\end{outcomes}

\begin{prereq}
\chref{what}: the four functions, and the habit of asking which of them a
business performs itself.
\end{prereq}

\begin{keyterms}
business model $\bullet$ revenue model $\bullet$ value proposition $\bullet$
B2C $\bullet$ B2B $\bullet$ C2C $\bullet$ C2B $\bullet$ marketplace $\bullet$
commission $\bullet$ subscription $\bullet$ freemium $\bullet$ advertising
$\bullet$ affiliate $\bullet$ unit economics $\bullet$ contribution margin
$\bullet$ customer acquisition cost $\bullet$ lifetime value $\bullet$
network effect
\end{keyterms}

\section{Two Questions That Get Confused}

Almost every failed online venture answered one of these two questions well and
the other not at all.

\begin{definitionbox}[Business Model and Revenue Model]
A \term{business model} says \emph{what value you create and for whom}: who the
customer is, what problem is solved, and why they would choose you rather than
the alternative. It is the answer to ``why does this deserve to exist?''

\smallskip
A \term{revenue model} says \emph{how the money reaches you}: who pays, for
what, how often, and how much. It is the answer to ``who writes the cheque?''

\smallskip
They are independent. The same business model --- say, connecting people who
want a ride with people who will drive --- can be monetised by commission per
ride, by a driver subscription, by advertising, or by charging nothing and
selling the data. Those four are different companies with different software.
\end{definitionbox}

The confusion matters because the two fail differently and the failures look
alike from outside. A venture with no business model and a clear revenue model
builds something nobody wants and charges correctly for it. A venture with a
strong business model and no revenue model builds something people love and
cannot pay for it --- and it is usually the second that students build.

\section{Classifying by Who Is at Each End}

The oldest classification asks only who is transacting with whom. It is coarse,
and it is still the first question worth asking, because it predicts a
surprising amount about the software.

\rowcolors{2}{white}{lightbg}
\begin{longtable}{@{}L{1.4cm}L{5.2cm}L{8.2cm}@{}}
\toprule
\rowcolor{primary!15}
\textbf{Type} & \textbf{Who to whom} & \textbf{What it implies for the
software} \\
\midrule
\endhead
\textbf{B2C} & A business sells to a consumer & Many small orders, card and
wallet payment, a catalogue tuned for browsing, and heavy traffic at
unpredictable moments. The default assumption of almost every storefront
platform, including the one this course uses \\
\textbf{B2B} & A business sells to another business & Few large orders,
negotiated and customer-specific prices, credit terms rather than immediate
payment, purchase orders, approval chains, and quantity breaks. Almost
everything in \chref{checkout} changes, and most of \chref{discounts} \\
\textbf{C2C} & A consumer sells to another consumer & The operator holds no
stock and never owns the goods. The hard parts move to trust, dispute
resolution and holding money in escrow --- \chref{auctions} is about this \\
\textbf{C2B} & A consumer sells to a business & Freelance marketplaces, stock
photography, influencer deals. The ``product'' is a person's time or licence,
which breaks the stock model entirely \\
\bottomrule
\end{longtable}

\begin{conceptbox}[Why the B2B/B2C split is the one that bites]
A platform built for B2C assumes one price per product. B2B does not have one
price per product --- it has a price per product \emph{per customer}, often
negotiated, often with quantity breaks, sometimes expiring.

\smallskip
This is the most common way a student project, or a real small business, picks
the wrong tool. Everything looks fine until the first wholesale customer asks
for their agreed rate, and the answer turns out to require a different data
model rather than a setting. \chref{planning} makes you ask this question
before you choose, which is the only time asking it is cheap.
\end{conceptbox}

\section{The Revenue Models}

There are not very many, and most businesses combine two or three.

\dsfig{%
\begin{tikzpicture}[
  rm/.style={rounded corners=3pt, draw=#1, line width=1pt, fill=#1!8,
             text=#1!60!black, font=\bfseries\scriptsize, align=center,
             text width=2.7cm, minimum height=1.5cm},
  lbl/.style={font=\scriptsize\itshape, text=neutral, align=center,
              text width=2.9cm}
]
  \node[rm=secondary] (a) at (0,0)    {Sales\\{\scriptsize\mdseries margin on
       each item sold}};
  \node[rm=greenhl]   (b) at (3.3,0)  {Commission\\{\scriptsize\mdseries a cut
       of somebody else's sale}};
  \node[rm=purplehl]  (c) at (6.6,0)  {Subscription\\{\scriptsize\mdseries paid
       access, repeating}};
  \node[rm=goldhl]    (d) at (9.9,0)  {Advertising\\{\scriptsize\mdseries
       attention sold to a third party}};
  \node[rm=tealhl]    (e) at (13.2,0) {Affiliate\\{\scriptsize\mdseries paid for
       sending a buyer}};

  \node[lbl] at (0,-1.55)    {you hold stock};
  \node[lbl] at (3.3,-1.55)  {you hold none};
  \node[lbl] at (6.6,-1.55)  {predictable; churn is the risk};
  \node[lbl] at (9.9,-1.55)  {needs scale before it pays};
  \node[lbl] at (13.2,-1.55) {no customer of your own};
\end{tikzpicture}}%
{The five revenue models a shop can use, and the single fact about each that
most shapes the software. Most real businesses run two of them at once.}%
{revenuemodels}

\subsection{Sales margin}

You buy at one price and sell at another. The oldest model, and the one the
storefront software in this course is built around by default.

Its distinguishing feature, from a software point of view, is that \emph{you
hold stock}. Stock means a quantity that can be wrong, can be reserved, can go
negative under concurrency, and must be reconciled against physical reality.
\chref{catalogue} and \chref{orders} spend real effort on this.

\subsection{Commission}

Somebody else sells; you take a percentage for having made it possible. This is
the marketplace model, and it is attractive because you hold no stock and your
costs barely rise with volume.

What you take on instead is everything in \chref{auctions}: two parties who do
not trust each other, money that must be held until both are satisfied, and
disputes that someone has to adjudicate. You have also acquired a problem the
sales model does not have --- you must attract \emph{both} sides, and neither
arrives before the other.

\subsection{Subscription}

The customer pays regularly for continued access or delivery. Revenue becomes
predictable, which is why investors like it and why so much of the web has
moved to it.

The software consequence is recurring billing, which is substantially harder
than one-off payment: cards expire, payments fail silently, customers pause and
resume, and the price may change mid-contract. \chref{gateway} explains why
a stored card is a different kind of integration from a single charge.

\subsection{Advertising and affiliate}

Advertising sells the attention of your visitors to somebody else. It needs
large traffic before it pays anything at all, which makes it a poor first
revenue model and a common fantasy.

Affiliate revenue is the mirror image: you send a visitor to somebody else's
shop and are paid if they buy. Cheap to start, and it leaves you without a
customer --- you never learn who bought, cannot serve them again, and the
relationship belongs to the merchant. \chref{seo} is where this becomes a
strategy rather than an accident.

\section{The Arithmetic of One Order}

A revenue model is a claim that can be checked, and the check is simpler than
students expect. Take one order and follow the money out of it.

\begin{worked}[Does this order make money?]
A shop sells a pair of shoes at \textbf{Rs.~4{,}500}, delivered anywhere in
Pakistan, with free returns. Work out what is left.

\smallskip
\textbf{Step 1 --- what came in.}\\
Revenue per order: \textbf{Rs.~4{,}500}.

\smallskip
\textbf{Step 2 --- the cost of the thing itself.}\\
The shoes cost the shop Rs.~2{,}800 to buy.\\
Remaining: $4{,}500 - 2{,}800 = \textbf{Rs.~1{,}700}$ (a gross margin of 38\%).

\smallskip
\textbf{Step 3 --- the cost of taking the money.}\\
The gateway charges 2.5\% plus Rs.~15 per transaction:
$0.025 \times 4{,}500 + 15 = 112.5 + 15 = \text{Rs.~127.50}$.\\
Remaining: $1{,}700 - 127.50 = \textbf{Rs.~1{,}572.50}$.

\smallskip
\textbf{Step 4 --- the cost of delivery.}\\
Courier, Rs.~250, absorbed by the shop because delivery was advertised free.\\
Remaining: $1{,}572.50 - 250 = \textbf{Rs.~1{,}322.50}$.

\smallskip
\textbf{Step 5 --- returns, which are not optional.}\\
20\% of shoe orders come back. A return costs the courier fee again (Rs.~250)
and the item cannot always be resold; assume a further Rs.~200 of loss on
average. Expected cost per \emph{order placed}:
$0.20 \times (250 + 200) = \text{Rs.~90}$.\\
Remaining: $1{,}322.50 - 90 = \textbf{Rs.~1{,}232.50}$.

\smallskip
\textbf{Step 6 --- the cost of finding the customer.}\\
The shop spends Rs.~60{,}000 a month on advertising and gets 50 orders from it:
a \term{customer acquisition cost} of $60{,}000 / 50 = \text{Rs.~1{,}200}$.\\
Remaining: $1{,}232.50 - 1{,}200 = \textbf{Rs.~32.50}$.

\smallskip
\textbf{The answer.} The order that looked like a 38\% margin earns
\textbf{Rs.~32.50}, which is \textbf{0.7\%} of the sale price --- and that is
before rent, salaries, hosting or tax. One extra return, or a 3\% gateway
instead of 2.5\%, and it is a loss.

\smallskip
\textbf{What this is really telling you.} The two numbers that decided it were
\emph{returns} and \emph{acquisition cost}, and neither appears anywhere in the
product's price. A shop that measures only gross margin believes it is making
38\%. This is why \chref{dashboard} is careful about which numbers on a
dashboard are worth acting on.
\end{worked}

\begin{definitionbox}[Contribution Margin, CAC and LTV]
\term{Contribution margin} is what one order contributes after all the costs
that vary with that order --- goods, payment fees, delivery, expected returns.
It is what is left to pay the costs that do not vary.

\smallskip
\term{Customer acquisition cost} (CAC) is the marketing spend divided by the
customers it produced.

\smallskip
\term{Lifetime value} (LTV) is the total contribution margin a customer
produces before they stop buying.

\smallskip
The test a business must pass is $\text{LTV} > \text{CAC}$, by a comfortable
multiple. The worked example above passes only if the customer orders again ---
its CAC is nearly its whole first-order margin, which is survivable for a
business with repeat customers and fatal for one without.
\end{definitionbox}

\begin{alertbox}[Where the network effect comes from, and where it does not]
A \term{network effect} exists when the service gets better for each user as
more users join. A marketplace has one: more sellers make it more useful to
buyers, which attracts more sellers.

\smallskip
A shop selling its own goods has \emph{no} network effect. Your thousandth
customer makes the shop no better for your first. Volume may buy you better
purchase prices, which is a real advantage, but it is not a network effect and
it does not protect you the same way.

\smallskip
This matters because the two justify completely different strategies. Running
at a loss to grow is rational when a network effect will eventually make you
defensible, and is just losing money when it will not. Plenty of ventures have
adopted the first strategy while having the second business.
\end{alertbox}

\begin{pitfall}
\begin{tight}
  \item \textbf{Naming a business model and thinking you have named a revenue
        model.} ``We are a marketplace for tutors'' says nothing about who pays.
  \item \textbf{Quoting gross margin as if it were profit.} The worked example
        went from 38\% to 0.7\% without anything unusual happening.
  \item \textbf{Forgetting returns.} In clothing and footwear they run at 20--40\%
        and they are the difference between a viable shop and a dead one.
  \item \textbf{Advertising as a first revenue model.} It pays at scale and
        nothing before it, so a new venture choosing it has chosen to earn
        nothing for a long time.
  \item \textbf{Assuming the platform is model-agnostic.} It is not. A B2C
        storefront has one price per product, and discovering that during
        implementation is expensive --- see \chref{planning}.
  \item \textbf{Treating ``free delivery'' as marketing.} It is a cost of
        Rs.~250 per order that appears in no price list and is paid out of the
        margin.
\end{tight}
\end{pitfall}

\begin{lab}[The model behind the shop you will build]
Still no software --- \chref{stack} starts that. This laboratory produces the
one-page brief that \chref{planning} turns into a decision and
\chref{project} is eventually marked against, so it is worth real effort.

\smallskip
\textbf{Choose a shop to build.} It must be plausible, small, and local enough
that you can check your assumptions: a bookshop, a bakery, a tailor, a
spare-parts dealer, a plant nursery. Do not choose ``an Amazon competitor''.

\smallskip
Then produce, on one page:

\begin{tightnum}
  \item \textbf{The business model in two sentences.} Who is the customer, what
        problem is solved, and why you rather than the shop down the road or
        the marketplace they already use.
  \item \textbf{The classification}: B2C, B2B, C2C or C2B --- and if you are
        tempted to say ``both'', say which one the \emph{software} must assume,
        because it can only assume one.
  \item \textbf{The revenue model}, from \dsref{revenuemodels}, with the reason
        for choosing it over the next most plausible one.
  \item \textbf{The arithmetic of one order}, following the six steps of the
        worked example with your own figures. Find real numbers where you can:
        gateway fees and courier rates are published, and a shopkeeper will
        usually tell you their return rate if you ask politely. Mark every
        number you guessed.
  \item \textbf{The two numbers that decide it.} From your own arithmetic, name
        the two inputs the result is most sensitive to, and say what would have
        to be true for the shop to lose money.
\end{tightnum}

\smallskip
\textbf{What to notice.} Almost everyone finds that the contribution margin is
far thinner than they expected, and that the deciding numbers are not the ones
they thought they were building a business around. That discovery is the point
of the exercise, and it is much cheaper to make here than in Chapter~23.
\end{lab}

\begin{checkpoint}
\begin{tightnum}
  \item A food delivery app charges restaurants 20\% commission, charges
        customers a delivery fee, and sells promoted placement in search
        results. Name the business model and all the revenue models, and say
        which single one you would expect to be most profitable per unit.
  \item A shop has a contribution margin of Rs.~400 per order and a customer
        acquisition cost of Rs.~1{,}500. Under what condition is this business
        viable? State it as an inequality and say what would need measuring.
  \item Explain why a B2B shop cannot usually be built by configuring a B2C
        platform, naming the specific assumption that breaks.
\end{tightnum}
\end{checkpoint}

\begin{chaptersummary}
\begin{tight}
  \item A business model says what value is created and for whom; a revenue
        model says who pays. They are independent, and a venture can have a
        good one and a fatal other.
  \item B2C, B2B, C2C and C2B classify by who transacts. The split that most
        affects the software is B2B against B2C: B2B has a price per customer,
        not a price per product.
  \item The five revenue models are sales margin, commission, subscription,
        advertising and affiliate. Holding stock or not is the fact about each
        that most shapes the system you have to build.
  \item Follow one order through six deductions --- goods, payment fee,
        delivery, returns, acquisition --- and the margin usually collapses. In
        the worked example, 38\% became 0.7\%.
  \item The viability test is $\text{LTV} > \text{CAC}$ by a comfortable
        multiple, which for a thin first-order margin means repeat custom is
        not optional.
  \item Marketplaces have network effects; shops selling their own goods do
        not, and strategies that assume one are expensive when it is absent.
\end{tight}
\end{chaptersummary}

\begin{reviewq}
  \item \textbf{(MCQ)} ``We connect home cooks with office workers who want
        lunch.'' This is: (a)~a revenue model; (b)~a business model; (c)~both;
        (d)~neither.
  \item \textbf{(MCQ)} Which revenue model most changes the \emph{payment}
        integration a shop needs? (a)~Sales margin; (b)~affiliate;
        (c)~subscription; (d)~advertising.
  \item \textbf{(Short)} Distinguish gross margin from contribution margin, and
        give two costs that appear in one and not the other.
  \item \textbf{(Short)} Explain why a marketplace has a harder launch than a
        shop, despite holding no stock.
  \item \textbf{(Short)} A clothing shop reports 45\% gross margin and is losing
        money. Give the two most likely explanations, and say what you would
        measure first.
  \item \textbf{(Applied)} A shop sells a product at Rs.~2{,}000 that costs it
        Rs.~1{,}100. Gateway: 2.9\% + Rs.~20. Delivery: Rs.~200, charged to the
        customer at Rs.~150. Returns: 12\%, each costing Rs.~200 in courier and
        Rs.~150 in unsellable stock. Compute the contribution margin per order,
        showing each step. Then find the return rate at which the shop breaks
        even.
  \item \textbf{(Applied)} Take the venture from this chapter's laboratory and
        argue for a \emph{different} revenue model than the one you chose. Give
        the strongest case you can, then say which software decisions in Parts
        III and IV would have to change if you adopted it.
\end{reviewq}
